\section{Reinforcement Learning}

In this project, reinforcement learning (RL) is used to improve the behavior of connected and automated vehicles (CAVs) in mixed traffic. The main idea is not to replace the baseline controller, but to add an adaptive layer on top of it. This learned layer adjusts the time headway policy according to the current traffic situation. In this way, the controller keeps its physical interpretation while gaining the ability to react to complex traffic patterns that are difficult to model explicitly.

\subsection{Fundamentals}

The RL problem is formulated as a Markov Decision Process (MDP),
\[
\mathcal{M} = (\mathcal{S}, \mathcal{A}, P, r),
\]
where $\mathcal{S}$ is the state space, $\mathcal{A}$ is the action space, $P$ represents the traffic dynamics, and $r$ is the reward.

At each time step, the agent observes the traffic state, selects an action, applies it to the controller, and receives a reward. In the current implementation, the state of each CAV $i$ is defined as
\[
x_i =
(\mu_{v,i}, \sigma_{v,i}^2, \delta^{\mathrm{macro}}_{v,i}, v_i, s_i, \Delta v_i, a_{\ell(i)}, \alpha_i^{prev}),
\]
where $\mu_{v,i}$ and $\sigma_{v,i}^2$ are the upstream mean velocity and velocity variance, $\delta^{\mathrm{macro}}_{v,i}=v_i-\mu_{v,i}$ is the speed mismatch with the local upstream flow, $v_i$ is the ego speed, $s_i$ is the spacing to the leader, $\Delta v_i$ is the relative speed, $a_{\ell(i)}$ is the leader acceleration, and $\alpha_i^{prev}$ is the previously applied headway scaling factor.

This state combines \emph{microscopic} information, such as spacing and relative velocity, with \emph{mesoscopic} information, such as the average speed and speed variance of upstream traffic. This is important because local measurements alone may not be enough to detect the development of traffic waves.

The action is a residual correction of the headway scaling:
\[
\mathcal{A} = \{\Delta\alpha_i \in [-\Delta\alpha_{\max}, +\Delta\alpha_{\max}]\}.
\]
The final headway scaling factor is then given by
\[
\alpha_i(t) =
\mathrm{clip}\left(\alpha_i^{rule}(t) + \Delta\alpha_i(t)\right),
\]
where $\alpha_i^{rule}(t)$ is the baseline headway scaling provided by the rule-based mesoscopic module. Therefore, RL only refines the baseline decision instead of replacing it. The learned policy can be written as
\[
\pi : \mathcal{S} \rightarrow \mathcal{A},
\qquad \Delta\alpha_i(t) = \pi(x_i(t)).
\]

\subsection{Reward design}

The reward is designed to stay aligned with the engineering goals of the project rather than with a purely numerical score:
\[
r_t = r_{\mathrm{sf}} + r_{\mathrm{sp}} + r_{\mathrm{ef}} + r_{\mathrm{cf}} + r_{\mathrm{ss}} + r_{\sigma^2}.
\]

For the spacing-regulation term, the controller uses the desired gap
\[
s_i^{\mathrm{des}}[t] = d_0 + h_i'[t]\,v_i[t],
\]
and the tracking error
\[
e_i[t] = s_i[t] - s_i^{\mathrm{des}}[t].
\]

The different reward terms have a clear physical meaning:
\begin{itemize}
    \item $r_{\mathrm{sf}}$: safety term, penalizing unsafe spacing and too small time headway;
    \item $r_{\mathrm{sp}}$: spacing regulation term, encouraging tracking of the desired spacing;
    \item $r_{\mathrm{ef}}$: efficiency term, encouraging the vehicle speed to remain close to the upstream mean speed;
    \item $r_{\mathrm{cf}}$: comfort term, penalizing high jerk;
    \item $r_{\mathrm{ss}}$: string stability term, penalizing amplification of spacing errors along the platoon;
    \item $r_{\sigma^2}$: traffic smoothness term, penalizing large velocity variance in the traffic stream.
\end{itemize}

Taken together, these terms encourage a policy that is safe first, but also smooth and useful at traffic scale. A controller that improved one vehicle while making the overall flow worse would not be acceptable here.

\subsection{PPO}

Proximal Policy Optimization (PPO) is a policy-gradient RL algorithm that updates the policy in a stable and conservative way. Its main advantage is that it is simple to implement and usually robust during training. PPO avoids large policy updates by using a clipped objective:
\[
L^{\mathrm{PPO}}(\theta)
=
\mathbb{E}_t
\left[
\min\left(
r_t(\theta)\hat{A}_t,\;
\mathrm{clip}(r_t(\theta),1-\epsilon,1+\epsilon)\hat{A}_t
\right)
\right],
\]
where
\[
r_t(\theta)=\frac{\pi_\theta(a_t|s_t)}{\pi_{\theta_{\mathrm{old}}}(a_t|s_t)},
\]
and $\hat{A}_t$ is the estimated advantage. In this project, PPO mainly serves as a strong baseline: it is simple to implement, reasonably stable, and easy to interpret when comparing learning behavior across runs.

\subsection{SAC}

Soft Actor-Critic (SAC) is another RL algorithm for continuous control. Unlike PPO, SAC is an off-policy method and includes an entropy term in its objective:
\[
J(\pi)
=
\sum_t
\mathbb{E}
\left[
r(s_t,a_t) + \beta \mathcal{H}(\pi(\cdot|s_t))
\right],
\]
where $\mathcal{H}(\pi(\cdot|s_t))$ is the policy entropy and $\beta$ controls the exploration level.

SAC is well suited to continuous-action problems and stochastic environments. Its replay buffer allows reuse of past experience, which can improve sample efficiency compared with on-policy methods such as PPO. This makes it a natural alternative here because the action is continuous and the mixed-traffic environment is noisy due to stochastic HDV behavior.

\subsection{Role of RL in the framework}

In our setup, RL does not produce the full acceleration command. It only adjusts the headway scaling around a rule-based mesoscopic baseline. This design choice is deliberate: it keeps the controller interpretable, limits unsafe departures from a known operating point, and makes it easier to judge whether learning is adding value beyond the handcrafted adaptation.
